\section{Physics CFD: Full-Device Stokes Model}

\subsection{Scientific Role}

The full-device CFD module provides a spatially resolved, continuum-flow
description of the asymmetric loop device.  The simulated domain contains the
incoming trunk, the upper bifurcation, the long left arm, the short right arm,
the lower merge junction, and the downstream outlet.  The central solid island
is treated as an internal no-slip boundary.

The purpose of this module is to generate a reusable background velocity field
for downstream physics enrichment and eventual machine-learning features.  The
current model is a single-phase, steady, incompressible Stokes approximation of
the continuous phase.  Droplets are not resolved as moving deformable bodies in
this CFD layer; their influence is represented elsewhere through occupancy and
hydraulic-state modules.

\[
\mathrm{centerline~geometry}
\rightarrow
\mathrm{Boolean~CAD~fluid~domain}
\rightarrow
\mathrm{triangular~finite\text{-}element~mesh}
\rightarrow
\mathrm{steady~Stokes~solution}
\rightarrow
\mathbf{u}(x,y),p(x,y).
\]

\subsection{Implemented Files}

\begin{itemize}
  \item \texttt{src/physics/full\_device\_cfd/domain.py}: constructs the
        full-device CAD polygon from calibrated centerlines.
  \item \texttt{src/physics/full\_device\_cfd/mesh.py}: generates and labels
        the constrained triangular finite-element mesh.
  \item \texttt{src/physics/full\_device\_cfd/solver.py}: assembles and solves
        the mixed Stokes system on the full-device mesh.
  \item \texttt{src/physics/full\_device\_cfd/field\_verification.py}: compares
        the production-candidate field against the fine reference field.
  \item \texttt{scripts/physics/build\_full\_device\_mesh.py}: writes the
        full-device CAD, mesh, and geometry diagnostics.
  \item \texttt{scripts/physics/run\_full\_device\_convergence.py}: runs scalar
        mesh convergence for the \(\alpha=0\) case.
  \item \texttt{scripts/physics/verify\_full\_device\_field\_convergence.py}:
        performs the final \(24~\mu\mathrm{m}\) versus
        \(12~\mu\mathrm{m}\) vector-field comparison.
\end{itemize}

\subsection{Configuration}

The production full-device CFD configuration is:

\begin{verbatim}
configs/physics/full_device_cfd.yml
\end{verbatim}

The selected production mesh size is frozen at exactly

\[
h_{\mathrm{target}} = 24~\mu\mathrm{m},
\qquad
h_{\mathrm{boundary}} = 12~\mu\mathrm{m}.
\]

This value was selected after comparing the complete \(24~\mu\mathrm{m}\)
velocity vector field against a \(12~\mu\mathrm{m}\) reference field on a common
Cartesian evaluation grid.

\subsection{Geometry Construction}

The device geometry is reconstructed from the extracted centerlines and fixed
channel dimensions, not from a segmentation mask.  Coordinates are stored in
the device Cartesian frame in micrometers.  The channel width and height are
both \(100~\mu\mathrm{m}\).

The CAD construction is deliberately geometric:

\[
\mathrm{centerlines}
\rightarrow
\mathrm{offset~walls~at}~\pm 50~\mu\mathrm{m}
\rightarrow
\mathrm{Boolean~union~of~connected~channels}
\rightarrow
\mathrm{central~island~hole}
\rightarrow
\mathrm{boundary~rings}.
\]

This construction yields one connected fluid component with one internal solid
island.  The inlet and outlet channels are blended into the loop as part of a
single continuous CAD domain so that no artificial wall fragments remain inside
the bifurcation or merge regions.

\begin{figure}[H]
  \centering
  \includegraphics[width=0.84\linewidth,height=0.70\textheight,keepaspectratio]{../outputs/physics/full_device_cfd/mesh/full_device_polygon_with_centerlines.png}
  \caption{Full-device fluid polygon reconstructed from calibrated centerlines.
  The blue region is the continuous fluid domain and the white central island is
  the no-slip solid boundary.}
  \label{fig:full-device-polygon}
\end{figure}

\begin{figure}[H]
  \centering
  \includegraphics[width=0.48\linewidth,height=0.48\textheight,keepaspectratio]{../outputs/physics/full_device_cfd/mesh/inlet_junction_polygon_zoom.png}
  \includegraphics[width=0.48\linewidth,height=0.48\textheight,keepaspectratio]{../outputs/physics/full_device_cfd/mesh/outlet_junction_polygon_zoom.png}
  \caption{Zoomed CAD diagnostics at the inlet split and outlet merge.  These
  checks verify that the branch transitions are continuous and do not contain
  internal wall remnants.}
  \label{fig:full-device-junction-cad}
\end{figure}

\subsection{Mesh Generation}

The mesher receives the final CAD polygon as a planar straight-line graph.  The
outer wall and inner island are sampled at the configured boundary spacing, and
a constrained triangular mesh is generated over the fluid region.  Boundary
facets are labeled as inlet, outlet, or wall for finite-element boundary
conditions.

For the frozen \(24~\mu\mathrm{m}\) production mesh:

\begin{center}
\begin{tabular}{lr}
\toprule
Quantity & Value \\
\midrule
Nodes & 2339 \\
Elements & 3807 \\
Minimum angle & \(25.727^\circ\) \\
Maximum angle & \(122.392^\circ\) \\
Median aspect ratio & 1.456 \\
Maximum aspect ratio & 2.285 \\
Fluid area & \(5.226\times10^5~\mu\mathrm{m}^2\) \\
Boundary length & \(1.043\times10^4~\mu\mathrm{m}\) \\
\bottomrule
\end{tabular}
\end{center}

\begin{figure}[H]
  \centering
  \includegraphics[width=0.84\linewidth,height=0.70\textheight,keepaspectratio]{../outputs/physics/full_device_cfd/mesh/full_device_mesh.png}
  \caption{Frozen \(24~\mu\mathrm{m}\) full-device triangular mesh.}
  \label{fig:full-device-mesh}
\end{figure}

\begin{figure}[H]
  \centering
  \includegraphics[width=0.84\linewidth,height=0.70\textheight,keepaspectratio]{../outputs/physics/full_device_cfd/mesh/full_device_mesh_quality.png}
  \caption{Mesh aspect-ratio diagnostic for the full-device production mesh.}
  \label{fig:full-device-mesh-quality}
\end{figure}

\begin{figure}[H]
  \centering
  \includegraphics[width=0.48\linewidth,height=0.48\textheight,keepaspectratio]{../outputs/physics/full_device_cfd/mesh/inlet_junction_mesh_zoom.png}
  \includegraphics[width=0.48\linewidth,height=0.48\textheight,keepaspectratio]{../outputs/physics/full_device_cfd/mesh/outlet_junction_mesh_zoom.png}
  \caption{Production mesh near the inlet split and outlet merge.}
  \label{fig:full-device-mesh-zooms}
\end{figure}

\subsection{Governing Equations}

The full-device solver uses the steady incompressible Stokes equations in the
two-dimensional channel midplane:

\[
-\mu\nabla^2\mathbf{u} + \nabla p + \alpha(\mathbf{x})\mathbf{u} = 0,
\qquad
\nabla\cdot\mathbf{u}=0.
\]

Here \(\mathbf{u}=(u_x,u_y)\) is the in-plane velocity, \(p\) is pressure,
\(\mu\) is dynamic viscosity, and \(\alpha(\mathbf{x})\) is an optional
distributed branch-resistance field.  For the validated \(\alpha=0\) baseline,

\[
\alpha_L=\alpha_R=0,
\]

so the equation reduces to the ordinary Stokes problem.  The solved
two-dimensional velocity is interpreted as a height-averaged background field
for the \(100~\mu\mathrm{m}\)-high channel.

\subsection{Boundary Conditions}

For the \(\alpha=0\) full-device validation case, the boundary formulation is:

\begin{itemize}
  \item prescribed parabolic inlet velocity profile with the experimental total
        flow rate;
  \item no-slip velocity on all channel walls and on the central island;
  \item natural traction condition at the downstream outlet;
  \item one pressure gauge degree of freedom to remove the constant pressure
        null mode;
  \item no pressure Dirichlet condition on the inlet or outlet.
\end{itemize}

The imposed inlet flux per unit depth is computed from the experimental total
volumetric flow rate \(Q\) and channel height \(H\):

\[
q_{\mathrm{in}} = \frac{Q}{H}.
\]

The mean inlet velocity follows from the channel width \(W\):

\[
U_{\mathrm{mean}} = \frac{q_{\mathrm{in}}}{W}.
\]

The prescribed planar Poiseuille profile is

\[
\mathbf{u}_{\mathrm{in}}(s)
=
\frac{3}{2}U_{\mathrm{mean}}
\left[1-\left(\frac{2s}{W}\right)^2\right]\mathbf{t}_{\mathrm{in}},
\qquad
-\frac{W}{2}\le s\le \frac{W}{2},
\]

where \(s\) is the transverse coordinate and \(\mathbf{t}_{\mathrm{in}}\) is
the inlet centerline tangent.

\subsection{Finite-Element Discretization}

The solver uses a Taylor--Hood mixed finite-element pair:

\[
\mathbf{u}_h \in [P_2]^2,
\qquad
p_h \in P_1.
\]

The weak form is assembled with scikit-fem.  Velocity Dirichlet conditions are
applied using labeled boundary facets so that P2 edge-midpoint velocity degrees
of freedom are retained and constrained correctly.  The linear saddle-point
system is solved with the direct backend.  No pressure-mass stabilization,
least-squares fallback, or artificial compressibility term is used.

\subsection{Validated \(\alpha=0\) Solution}

The first validated full-device solution is saved under:

\begin{verbatim}
outputs/physics/full_device_cfd/alpha0_equal_split_smoke/
\end{verbatim}

The historical frozen baseline snapshot is:

\begin{verbatim}
outputs/physics/full_device_cfd/baseline_alpha0_v1/
\end{verbatim}

The \(\alpha=0\) case does not produce a 50/50 split in the asymmetric loop,
because the long and short branches have different hydraulic resistances.  The
computed branch split is:

\[
f_L = \frac{Q_L}{Q_L+Q_R}=0.458119,
\qquad
f_R=0.541881.
\]

Representative scalar quantities for the \(24~\mu\mathrm{m}\) mesh are:

\begin{center}
\begin{tabular}{lr}
\toprule
Quantity & Value \\
\midrule
Solver backend & scikit-fem/direct \\
Inlet flux & \(-5.4444\times10^{-6}~\mathrm{m^2/s}\) \\
Outlet flux & \(5.4774\times10^{-6}~\mathrm{m^2/s}\) \\
Left branch flux & \(2.4932\times10^{-6}~\mathrm{m^2/s}\) \\
Right branch flux & \(2.9490\times10^{-6}~\mathrm{m^2/s}\) \\
Maximum velocity & \(8.1675\times10^{-2}~\mathrm{m/s}\) \\
Pressure range & \(166.436~\mathrm{Pa}\) \\
\bottomrule
\end{tabular}
\end{center}

\begin{figure}[H]
  \centering
  \includegraphics[width=0.84\linewidth,height=0.70\textheight,keepaspectratio]{../outputs/physics/full_device_cfd/alpha0_equal_split_smoke/diagnostics/speed_magnitude.png}
  \caption{Velocity magnitude for the validated full-device \(\alpha=0\)
  Stokes solution.}
  \label{fig:full-device-speed}
\end{figure}

\begin{figure}[H]
  \centering
  \includegraphics[width=0.84\linewidth,height=0.70\textheight,keepaspectratio]{../outputs/physics/full_device_cfd/alpha0_equal_split_smoke/diagnostics/pressure.png}
  \caption{Pressure field for the validated full-device \(\alpha=0\) Stokes
  solution.}
  \label{fig:full-device-pressure}
\end{figure}

\begin{figure}[H]
  \centering
  \includegraphics[width=0.84\linewidth,height=0.70\textheight,keepaspectratio]{../outputs/physics/full_device_cfd/alpha0_equal_split_smoke/diagnostics/full_device_streamlines_very_dense_postprocessed.png}
  \caption{Dense post-processed streamlines sampled from the solved
  full-device velocity field.  These streamlines are visualization samples of a
  continuous FEM velocity field, not a limit on CFD resolution.}
  \label{fig:full-device-streamlines}
\end{figure}

\subsection{Scalar Mesh Convergence}

A scalar convergence study was run for target element sizes
\(48,36,24,18,12~\mu\mathrm{m}\).  The \(12~\mu\mathrm{m}\) case was used as
the fine reference.  The branch split and maximum velocity are already stable
at \(24~\mu\mathrm{m}\), while the pressure range continues to move by less
than \(0.5\%\) relative to the \(12~\mu\mathrm{m}\) reference.

\begin{center}
\begin{tabular}{rrrrrr}
\toprule
\(h\) (\(\mu\mathrm{m}\)) & Elements & \(f_L\) & \(|\Delta f_L|\) vs 12 & \(v_{\max}\) rel. error & Pressure range rel. error \\
\midrule
48 & 1040 & 0.458190 & \(3.36\times10^{-6}\) & \(4.32\times10^{-4}\) & \(1.50\times10^{-2}\) \\
36 & 1768 & 0.456614 & \(1.58\times10^{-3}\) & \(1.26\times10^{-3}\) & \(1.28\times10^{-2}\) \\
24 & 3807 & 0.458119 & \(7.41\times10^{-5}\) & \(8.59\times10^{-4}\) & \(4.87\times10^{-3}\) \\
18 & 6508 & 0.458378 & \(1.85\times10^{-4}\) & \(5.82\times10^{-4}\) & \(2.34\times10^{-3}\) \\
12 & 14154 & 0.458193 & reference & reference & reference \\
\bottomrule
\end{tabular}
\end{center}

\begin{figure}[H]
  \centering
  \includegraphics[width=0.95\linewidth,height=0.55\textheight,keepaspectratio]{../outputs/physics/full_device_cfd/convergence_alpha0_fine/full_device_alpha0_convergence.png}
  \caption{Scalar mesh convergence for the full-device \(\alpha=0\) Stokes
  solution.}
  \label{fig:full-device-scalar-convergence}
\end{figure}

\subsection{Field-Level Mesh Verification}

The production decision was made using a field-level comparison between the
\(24~\mu\mathrm{m}\) candidate and the \(12~\mu\mathrm{m}\) reference.  Both
solutions were evaluated on the same Cartesian grid with
\(6~\mu\mathrm{m}\) spacing.  Points outside the final fluid polygon and points
inside the central island were masked.  Velocity components were interpolated
from the FEM basis directly and compared in SI units.

For each region, the vector-field error was computed as

\[
E_{\mathbf{u}}
=
\frac{\|\mathbf{u}_{24}-\mathbf{u}_{12}\|_2}
     {\|\mathbf{u}_{12}\|_2}.
\]

Angular error was computed from

\[
\theta =
\cos^{-1}
\left[
\mathrm{clip}
\left(
\frac{\mathbf{u}_{24}\cdot\mathbf{u}_{12}}
     {|\mathbf{u}_{24}||\mathbf{u}_{12}|},
-1,1
\right)
\right],
\]

excluding points with reference speed below \(1\%\) of the maximum reference
speed.  Relative speed errors used the same \(1\%\) speed floor to avoid
letting stagnation-point noise dominate the decision.

\begin{center}
\begin{tabular}{lrrrr}
\toprule
Region & \(E_{\mathbf{u}}\) & Median angle & 95th angle & Speed L2 error \\
\midrule
Full domain & 0.009808 & \(0.0055^\circ\) & \(0.336^\circ\) & 0.008221 \\
Inlet split junction & 0.012043 & \(0.0089^\circ\) & \(0.829^\circ\) & 0.009419 \\
Outlet merge junction & 0.011063 & \(0.0127^\circ\) & \(0.865^\circ\) & 0.008380 \\
\bottomrule
\end{tabular}
\end{center}

The inlet separatrix was estimated by tracing the same inlet seed set in both
fields.  The transition seed location differed by

\[
1.533~\mu\mathrm{m}
= 1.53\%~\mathrm{of~the~channel~width}.
\]

All field-level acceptance criteria passed:

\begin{itemize}
  \item full-domain vector relative L2 error below \(2\%\);
  \item inlet-junction vector relative L2 error below \(2\%\);
  \item outlet-junction vector relative L2 error below \(2\%\);
  \item median angular error below \(1^\circ\);
  \item 95th-percentile angular error below \(3^\circ\);
  \item separatrix-location difference below \(2\%\) of channel width.
\end{itemize}

\begin{figure}[H]
  \centering
  \includegraphics[width=0.48\linewidth,height=0.48\textheight,keepaspectratio]{../outputs/physics/full_device_cfd/convergence/field_comparison_24_vs_12/velocity_magnitude_24um.png}
  \includegraphics[width=0.48\linewidth,height=0.48\textheight,keepaspectratio]{../outputs/physics/full_device_cfd/convergence/field_comparison_24_vs_12/velocity_magnitude_12um.png}
  \caption{Velocity magnitude on the common grid for the \(24~\mu\mathrm{m}\)
  candidate and \(12~\mu\mathrm{m}\) reference.  Both panels use identical
  color limits.}
  \label{fig:field-verification-magnitudes}
\end{figure}

\begin{figure}[H]
  \centering
  \includegraphics[width=0.48\linewidth,height=0.48\textheight,keepaspectratio]{../outputs/physics/full_device_cfd/convergence/field_comparison_24_vs_12/absolute_vector_error_24_vs_12.png}
  \includegraphics[width=0.48\linewidth,height=0.48\textheight,keepaspectratio]{../outputs/physics/full_device_cfd/convergence/field_comparison_24_vs_12/angular_error_24_vs_12.png}
  \caption{Absolute vector error and angular error between the \(24~\mu\mathrm{m}\)
  and \(12~\mu\mathrm{m}\) full-device velocity fields.}
  \label{fig:field-verification-errors}
\end{figure}

\begin{figure}[H]
  \centering
  \includegraphics[width=0.48\linewidth,height=0.48\textheight,keepaspectratio]{../outputs/physics/full_device_cfd/convergence/field_comparison_24_vs_12/inlet_junction_error_zoom.png}
  \includegraphics[width=0.48\linewidth,height=0.48\textheight,keepaspectratio]{../outputs/physics/full_device_cfd/convergence/field_comparison_24_vs_12/outlet_junction_error_zoom.png}
  \caption{Field-level error zooms at the full-device inlet split and outlet
  merge. Both fields come from the complete loop geometry.}
  \label{fig:field-verification-junction-zooms}
\end{figure}

\subsection{Flux Diagnostic}

The scalar convergence study showed a non-monotone inlet--outlet flux mismatch.
The flux audit confirmed that the current reported boundary fluxes are
computed from boundary facets using the FEM P2 velocity field, not from a
coarse plotted grid.  Each boundary facet is evaluated at its two endpoints and
midpoint with Simpson weights.

For the field-verification runs:

\begin{center}
\begin{tabular}{lrr}
\toprule
Case & Inlet/outlet mismatch & Relative mismatch \\
\midrule
\(24~\mu\mathrm{m}\) & \(3.298\times10^{-8}~\mathrm{m^2/s}\) & 0.00606 \\
\(12~\mu\mathrm{m}\) & \(1.541\times10^{-8}~\mathrm{m^2/s}\) & 0.00283 \\
\bottomrule
\end{tabular}
\end{center}

This scalar mismatch is documented but is not used to reject the
\(24~\mu\mathrm{m}\) production mesh, because the field-level errors, branch
split, pressure range, and separatrix comparison all satisfy the production
criteria.

\subsection{Production Decision}

The full-device production mesh is frozen at:

\[
\boxed{h_{\mathrm{target}}=24~\mu\mathrm{m}.}
\]

The final recommendation is to use the \(24~\mu\mathrm{m}\) mesh as the default
full-device CFD mesh.  The \(12~\mu\mathrm{m}\) mesh remains the fine reference
for verification, not the default production mesh.

\subsection{Execution}

Build the full-device CAD and production mesh:

\begin{verbatim}
.\.venv\Scripts\python.exe -m scripts.physics.build_full_device_mesh
\end{verbatim}

Run the scalar mesh convergence study:

\begin{verbatim}
$env:MPLCONFIGDIR='outputs\.matplotlib'
.\.venv\Scripts\python.exe scripts\physics\run_full_device_convergence.py `
  --target-sizes-um 48 36 24 18 12 `
  --output-dir outputs\physics\full_device_cfd\convergence_alpha0_fine
\end{verbatim}

Run the final field-level verification:

\begin{verbatim}
$env:MPLCONFIGDIR='outputs\.matplotlib'
.\.venv\Scripts\python.exe scripts\physics\verify_full_device_field_convergence.py `
  --grid-spacing-um 6 `
  --output-dir outputs\physics\full_device_cfd\convergence\field_comparison_24_vs_12
\end{verbatim}

The field-verification artifacts are saved in:

\begin{verbatim}
outputs/physics/full_device_cfd/convergence/field_comparison_24_vs_12/
\end{verbatim}
